\section{Reproducibility Checklist}
\label{app:reproducibility}

\begin{itemize}[leftmargin=*,itemsep=2pt]
  \item Freeze benchmark commit, task list, task order, evaluator, and database seed.
  \item Record exact provider-returned model identifier, non-thinking mode,
        temperature, prompt hash, schema hash, and API Predictor output.
  \item Use unique experiment, task, trial, and configuration identifiers.
  \item Persist the final environment state before invoking the evaluator.
  \item Save every model proposal, runtime disposition (execute, defer, reject),
        tool result, state reduction, verification, recovery, and completion decision.
  \item Close a deferred \texttt{tool\_call\_id} with one
        \texttt{deferred/not\_executed} result; any later proposal uses a new ID.
  \item Substitute dependent values only for explicit Action-IR placeholders with
        a declared edge and one unambiguous observed output; never repair guessed IDs.
  \item Separate agent, predictor, user-simulator, judge, and runtime token usage.
  \item Classify infrastructure failures and rerun only under the identical frozen condition.
  \item Redact API keys, credentials, personal information, and sensitive tool fields from traces.
  \item Run the result-macro validator before compiling the paper; projected-target
        macros are forbidden in manuscript sources.
  \item Report the code commit, environment lock hash, source-data hashes, and
        generated-result hashes with the final artifact.
\end{itemize}

\subsection{Frozen diagnostic configuration}

\begin{center}
\small
\begin{tabularx}{\linewidth}{@{}lYY@{}}
\toprule
Field & AppWorld diagnostic & $\tau^2$-bench diagnostic \\
\midrule
Agent model ID & \AWModelID & \TauAgentModelID \\
Endpoint & \AWAPIEndpoint & \TauAgentEndpoint \\
Access date & \AWModelAccessDate & \TauModelAccessDate \\
Provider revision & \AWProviderVersion & \TauProviderVersion \\
Benchmark package/revision & AppWorld revision not archived & $\tau^2$ \TauPackageVersion; \texttt{1901a301961c} (full hash in manifest) \\
Thinking & disabled & disabled \\
Temperature & \ModelTemperature & \ModelTemperature \\
\bottomrule
\end{tabularx}
\end{center}

The AppWorld runner used an API Predictor capped at \AWAPILimit{} active APIs and
an \AWMaxTurns{}-turn cap. Provider-returned immutable build identifiers were not
preserved in the historical runs, so the endpoint, model alias, and access dates
bound the result more precisely than an unsupported version claim. The historical
run is exploratory because its task-level traces informed development. Formal reruns
must persist the response model field, system fingerprint when available, and
deployment revision when supplied.

\subsection{Runner repair audit}

The historical runner repair log records a state-save correction and
\AWGBKRepairCount{} Windows GBK decoding failures rerun under interpreter-level
UTF-8, of which \AWGBKRepairSuccessCount{} subsequently passed. The original
runtime package version, environment lock, and repair execution logs were not
archived together, so these counts are not used as a performance result or a
reproducible repair-rate claim. Every future condition must record the persistence
event, environment lock, runner revision, and exact rerun reason; model failures
must never be selectively rerun. The available historical identifiers are:

\begingroup
\raggedright
\begin{itemize}[leftmargin=*,itemsep=2pt]
  \item canonical sorted AppWorld task list SHA-256:
        \texttt{\AWTaskListHashDisplay};
  \item UTF-8 repair subset SHA-256:
        \texttt{\AWRepairListHashDisplay}.
\end{itemize}
\endgroup

\AWInvalidDurationCount{} historical negative duration values produced by environment time freezing
are excluded from the archived duration summaries only. Future conditions use a
monotonic host clock and retain the raw timing stream.

\section{Implementation Inventory}
\label{app:implementation}

For artifact inspection only, the current workspace contains a framework-neutral
runtime interface, versioned contracts, a DeerFlow hook, one tested AppWorld
integration, and prototype $\tau^2$-bench adapters. Code size and schema counts
are intentionally excluded from the main argument: they describe artifact scope,
not correctness. The current workspace has
no recoverable pre-exposure Git history, so it cannot prove a historical design
freeze. The public release must publish a source revision, dependency lock, and
CI log before confirmatory benchmark access.

\begin{center}
\small
\begin{tabularx}{\linewidth}{@{}lY@{}}
\toprule
Scope & Local validation outcome and boundary \\
\midrule
Runtime and AppWorld integration & \texttt{python -m pytest} over configured paths: \LocalUnitTestPassed{} passed. Contract, state-store, guard, trace, and public-catalog tests passed locally. \\
End-to-end runtime prototype & \PrototypeValidationPassed{}/\PrototypeValidationRepetitionCount{} repetitions of one deterministic mock regression scenario passed; \PrototypeValidationCrashes{} crashes. Each used SQLite state and trace stores, rejected one invalid call, verified one write, blocked unverified completion, prevented a duplicate write, then reloaded durable state and trace. This is not an AppWorld score. \\
DeerFlow live compatibility & \LocalUnitTestSkipped{} skipped because \texttt{deerflow.sandbox.tools} is not installed. This integration is not verified. \\
AppWorld public catalog & The real API-document fixture yields \AppWorldCatalogToolCount{} tools. Schema conversion is unit-tested; this is not an end-to-end benchmark run. \\
$\tau^2$ adapter & Local smoke invocation only. No repeated-trial or official Pass$^k$ claim is made. \\
\bottomrule
\end{tabularx}
\end{center}

The test outcome above is a local validation record, not a substitute for an
external CI run, baseline-bypass equivalence test, or evaluator-confirmed agent
success. Those release checks remain required before a confirmatory-results paper.
